\section*{第 \ref{ch:g12:curly-arrows} 章——反应机理：弯箭头}

\begin{solution}{exo:g12:curly-arrows:1}
\ce{OH-}：\omterm{def:g12:curly-arrows:site}{给电子位点}（孤对\omterm{def:g9:inside-the-atom:nucleus}{电子}，带负电荷）；无\omterm{def:g12:curly-arrows:site}{受电子位点}。
\ce{H+}：\omterm{def:g12:curly-arrows:site}{受电子位点}（完全没有\omterm{def:g9:inside-the-atom:nucleus}{电子}）。\ce{CH2=CH2}：\omterm{def:g12:curly-arrows:site}{给电子位点}（\omterm{def:g10:lewis-and-shape:multiple-bond}{双
键}）。\ce{CH3-Cl}：\omterm{def:g12:curly-arrows:site}{给电子位点}（\ce{Cl} 的孤对\omterm{def:g9:inside-the-atom:nucleus}{电子}），\omterm{def:g12:curly-arrows:site}{受电子位点}（碳，
$\delta^+$）。\ce{H2O}：\omterm{def:g12:curly-arrows:site}{给电子位点}（\ce{O} 的孤对\omterm{def:g9:inside-the-atom:nucleus}{电子}）；其氢\omterm{def:g7:atoms-and-molecules:atom}{原子}
（$\delta^+$）是较弱的\omterm{def:g12:curly-arrows:site}{受电子位点}。
\end{solution}

\begin{solution}{exo:g12:curly-arrows:2}
(a) \omterm{def:g12:curly-arrows:categories}{取代}；(b) \omterm{def:g12:curly-arrows:categories}{加成}；(c) \omterm{def:g12:curly-arrows:categories}{消除}；(d) \omterm{def:g12:curly-arrows:categories}{加成}。
\end{solution}

\begin{solution}{exo:g12:curly-arrows:3}
箭头起始于一对\omterm{def:g9:inside-the-atom:nucleus}{电子}（孤对\omterm{def:g9:inside-the-atom:nucleus}{电子}或化学键），终点在
\omterm{def:g9:inside-the-atom:nucleus}{电子}对所去的\omterm{def:g7:atoms-and-molecules:atom}{原子}上。若由氧的孤对\omterm{def:g9:inside-the-atom:nucleus}{电子}指向碳，
则该\omterm{def:g9:inside-the-atom:nucleus}{电子}对形成新的 \ce{O-C} 键。
\end{solution}

\begin{solution}{exo:g12:curly-arrows:4}
氨提供氮\omterm{def:g7:atoms-and-molecules:atom}{原子}上的孤对\omterm{def:g9:inside-the-atom:nucleus}{电子}；箭头由该
孤对\omterm{def:g9:inside-the-atom:nucleus}{电子}指向 \ce{H+}。生成物 \ce{NH4+} 带电荷 $+1$。
\end{solution}

\begin{solution}{exo:g12:curly-arrows:5}
一步反应中生成、在后续某步中被消耗的物种；它不在
总反应方程式中出现。此处为碳正\omterm{def:g9:ions:ion}{离子} \ce{CH3-CH2+}。
\end{solution}

\begin{solution}{exo:g12:curly-arrows:6}
一个箭头从 \ce{OH-} 中氧的孤对\omterm{def:g9:inside-the-atom:nucleus}{电子}指向碳；
另一个箭头从 \ce{C-Br} 键指向溴。电荷：左侧为 $-1$
（\ce{OH-}）；右侧甲醇为中性，\ce{Br-} 带 $-1$：
两侧总电荷均为 $-1$。
\end{solution}

\begin{solution}{exo:g12:curly-arrows:7}
溴的\omterm{def:g11:polarity-and-cohesion:electronegativity}{电负性}比氢大：它取得该\omterm{def:g9:inside-the-atom:nucleus}{电子}对，
以 \ce{Br-}（$-1$）形式离去；氢与一个碳成键，
另一个碳则带正电荷：\ce{CH3-CH2+}（$+1$）。
\end{solution}

\begin{solution}{exo:g12:curly-arrows:8}
1.~\ce{CH3-CH2-OH + H+ -> CH3-CH2-OH2+}：\omterm{def:g12:curly-arrows:site}{给电子位点}为
氧的孤对\omterm{def:g9:inside-the-atom:nucleus}{电子}；\omterm{def:g12:curly-arrows:site}{受电子位点}为 \ce{H+}。2.~\ce{CH3-CH2-OH2+ -> CH3-CH2+ + H2O}：
\ce{C-O} \omterm{def:g9:inside-the-atom:nucleus}{电子}对转移给带正电的氧（\omterm{def:g12:curly-arrows:site}{受电子位点}），该氧随之
以水的形式离去。3.~\ce{CH3-CH2+ -> CH2=CH2 + H+}：\ce{C-H} 键的\omterm{def:g9:inside-the-atom:nucleus}{电子}对
（\omterm{def:g12:curly-arrows:site}{给电子位点}）移向带正电的碳（\omterm{def:g12:curly-arrows:site}{受电子位点}），形成\omterm{def:g10:lewis-and-shape:multiple-bond}{双键}，
\ce{H+} 离去。第 1 步中消耗的氢\omterm{def:g9:ions:ion}{离子}，
在第 3 步中被放回，它在方程两侧均出现而相互抵消。
\end{solution}

\begin{solution}{exo:g12:curly-arrows:9}
箭头从水中氧的孤对\omterm{def:g9:inside-the-atom:nucleus}{电子}指向带正电的碳。
失去 \ce{H+} 后生成乙醇 \ce{CH3-CH2-OH}。总反应为：
\ce{CH2=CH2 + H2O -> CH3-CH2-OH}，为\omterm{def:g12:curly-arrows:categories}{加成}反应。
\end{solution}

\begin{solution}{exo:g12:curly-arrows:10}
两对\omterm{def:g9:inside-the-atom:nucleus}{电子}在一步中同时移动：\ce{C-O} 键形成，\ce{C-Br}
键断裂。不生成\omterm{def:g12:curly-arrows:mechanism}{反应中间体}。
\end{solution}

\begin{solution}{exo:g12:curly-arrows:11}
氯的\omterm{def:g11:polarity-and-cohesion:electronegativity}{电负性}差（$3.16 - 2.55 = 0.61$）大于
溴的（$2.96 - 2.55 = 0.41$）：氯甲烷中的碳更
$\delta^+$，仅此一点就会使其反应更快。但事实并非如此：反应
速率还取决于\omterm{def:g10:electron-shells:family}{卤素}带着\omterm{def:g9:inside-the-atom:nucleus}{电子}对离去的难易程度，
而体积更大的溴更易离去。
\end{solution}

\begin{solution}{exo:g12:curly-arrows:12}
若氢与 1 号碳成键：\ce{CH3-CH+-CH3}，进而
得到 2-氯丙烷 \ce{CH3-CHCl-CH3}。若氢与 2 号碳成键：
\ce{+CH2-CH2-CH3}，则得到 1-氯丙烷 \ce{CH2Cl-CH2-CH3}。
\end{solution}

\begin{solution}{exo:g12:curly-arrows:13}
箭头方向画反了：它应从\omterm{def:g12:curly-arrows:site}{给电子位点}出发，即
\ce{OH-} 的一对孤对\omterm{def:g9:inside-the-atom:nucleus}{电子}，并指向\omterm{def:g12:curly-arrows:site}{受电子位点}碳。新键应
在氧与碳之间形成，而非氧与溴之间：溴离去时带走
\ce{C-Br} 键的\omterm{def:g9:inside-the-atom:nucleus}{电子}对。
\end{solution}

\begin{solution}{exo:g12:curly-arrows:14}
酸的 \ce{C=O} 基团中，碳与两个\omterm{def:g11:polarity-and-cohesion:electronegativity}{电负性}较强的
氧\omterm{def:g7:atoms-and-molecules:atom}{原子}相连，带较强的 $\delta^+$，是\omterm{def:g12:curly-arrows:site}{受电子位点}。醇中的
氧带 $\delta^-$ 且有两对孤对\omterm{def:g9:inside-the-atom:nucleus}{电子}，是\omterm{def:g12:curly-arrows:site}{给电子位点}。
总体失去一\omterm{def:g7:atoms-and-molecules:molecule}{分子}水；酸中的 \ce{-OH} 被醇的 \ce{-O-R} \omterm{def:g12:curly-arrows:categories}{取代}，
为\omterm{def:g12:curly-arrows:categories}{取代}反应。
\end{solution}

\begin{solution}{exo:g12:curly-arrows:15}
在\omterm{def:g10:lewis-and-shape:multiple-bond}{双键}\omterm{def:g9:inside-the-atom:nucleus}{电子}的排斥下，\ce{Br-Br} 键的\omterm{def:g9:inside-the-atom:nucleus}{电子}对
移向远端的溴，使近端的溴带 $\delta^+$，成为
\omterm{def:g12:curly-arrows:site}{受电子位点}。第一步：一个箭头由\omterm{def:g10:lewis-and-shape:multiple-bond}{双键}指向近端的溴，
另一个箭头由 \ce{Br-Br} 键指向远端的溴，后者以
\ce{Br-} 的形式离去；
原\omterm{def:g10:lewis-and-shape:multiple-bond}{双键}的另一个碳\omterm{def:g7:atoms-and-molecules:atom}{原子}上留下正电荷。
\end{solution}

\begin{solution}{pb:g12:curly-arrows:1}
\textbf{1.} 乙烯：\ce{H2C=CH2}，含一个\omterm{def:g10:lewis-and-shape:multiple-bond}{双键}和四个 \ce{C-H} 键。
溴化氢：\ce{H-Br}，溴\omterm{def:g7:atoms-and-molecules:atom}{原子}上有三对孤对\omterm{def:g9:inside-the-atom:nucleus}{电子}。

\textbf{2.} \omterm{def:g10:lewis-and-shape:multiple-bond}{双键}的第二对\omterm{def:g9:inside-the-atom:nucleus}{电子}
构成一个富\omterm{def:g9:inside-the-atom:nucleus}{电子}区域。

\textbf{3.} $2.96 - 2.20 = 0.76$：\ce{H} 带 $\delta^+$。

\textbf{4.} 氢\omterm{def:g7:atoms-and-molecules:atom}{原子}。

\textbf{5.} 一个箭头由\omterm{def:g10:lewis-and-shape:multiple-bond}{双键}指向 \ce{HBr} 中的氢；
另一个箭头由 \ce{H-Br} 键指向溴。

\textbf{6.} 碳正\omterm{def:g9:ions:ion}{离子} \ce{CH3-CH2+}（$+1$）和 \ce{Br-}（$-1$）。

\textbf{7.} 反应前：$0 + 0 = 0$；反应后：$+1 - 1 = 0$。

\textbf{8.} 一个箭头由 \ce{Br-} 的一对孤对\omterm{def:g9:inside-the-atom:nucleus}{电子}指向带正电的
碳\omterm{def:g7:atoms-and-molecules:atom}{原子}。

\textbf{9.} 六个\omterm{def:g9:inside-the-atom:nucleus}{电子}（三个键）：缺一对\omterm{def:g9:inside-the-atom:nucleus}{电子}即可达到八\omterm{def:g9:inside-the-atom:nucleus}{电子}结构，
因此它能接受附近任一供体提供的一对\omterm{def:g9:inside-the-atom:nucleus}{电子}。

\textbf{10.} 碳正\omterm{def:g9:ions:ion}{离子} \ce{CH3-CH2+}。

\textbf{11.} 它在第一步生成，在第二步被消耗。

\textbf{12.} \omterm{def:g12:curly-arrows:categories}{加成}反应。

\textbf{13.} 只有\omterm{def:g11:functional-groups:functional-group}{官能团}发生变化：两个碳组成的碳链保持不变，
\omterm{def:g10:lewis-and-shape:multiple-bond}{双键}变为单键，并分别连有 \ce{H} 和 \ce{Br}。

\textbf{14.} \qty{28.0}{g/mol}；\qty{108.9}{g/mol}。

\textbf{15.} $10.0 / 28.0 = \qty{0.357}{mol}$。

\textbf{16.} 乙烯：$0.357 - x$；\ce{HBr}：过量；溴乙烷：$x$。
乙烯为\omterm{def:g10:reaction-progress-table:limiting}{限量反应物}：$x_{\max} = \qty{0.357}{mol}$。

\textbf{17.} $0.357 \times 108.9 = \qty{38.9}{g}$。

\textbf{18.} $0.75 \times 38.9 = \qty{29.2}{g}$。

\textbf{19.} 约 \qty{29}{g} 溴乙烷。
\end{solution}
